\chapter{Warum die Kreuzsummenrelation Kürzbarkeit braucht}

Auf der zweielementigen Menge \(\{0,1\}\) liefert das Maximum ein
kommutatives Monoid mit neutralem Element \(0\). Die Operation ist jedoch
nicht kürzbar. Wir prüfen die Struktur und anschließend das Scheitern der
Transitivität der Kreuzsummenrelation. Alle Zeichen mit Index \(0\) werden
in diesem Kapitel konkret festgelegt; eine Kürzbarkeitsannahme wird dabei
nicht gemacht.

\par\Needspace{16\baselineskip}
\section{Der geordnete Träger}

\par\Needspace{12\baselineskip}
\hypertarget{differences.statement.FDMaxCarrier}{}
\FormulaDefDeltaK[Der zweielementige Träger]{%
  M_0\coloneqq\{0,1\}
}{FDMaxCarrier}{%
  \DeltaRow{Mengen}{M_0\dsep0\dsep1}
}

\par\Needspace{12\baselineskip}
\hypertarget{differences.statement.FDMaxNaturalConstants}{}
\FormulaThmDeltaK[Null und Eins sind verschiedene natürliche Zahlen]{%
  0\in\mathbb N\land1\in\mathbb N\land0\neq1
}{FDMaxNaturalConstants}{%
  \DeltaRow{Mengen}{0\dsep1}
}
\begin{tabproofwide}
  \proofstepwidestar[]{0\in\mathbb N}{\FormulaRefAuto{PeanoZero}[ax]}
  \proofstepwidestar[]{\Succ(0)\in\mathbb N}{\FormulaRefAuto{PeanoSuccClosure}[ax]{1}}
  \proofstepwidestar[]{1=\Succ(0)}{\FormulaRefAuto{PeanoOne}[def]}
  \proofstepwidestar[]{\Succ(0)=1}{\FormulaRefAuto{a=b\vdash b=a}[thm]{3}}
  \proofstepwidestar[]{1\in\mathbb N}{\rIE{4,2}}
  \proofstepwidestar[]{\Succ(0)\neq0}{\FormulaRefAuto{PeanoSuccNonzero}[ax]{1}}
  \proofstepwidestar[7]{0=1}{\rA}
  \proofstepwidestar[7]{0=\Succ(0)}{\rIE{3,7}}
  \proofstepwidestar[7]{\Succ(0)=0}{\FormulaRefAuto{a=b\vdash b=a}[thm]{8}}
  \proofstepwidestar[7]{\bot}{\rBI{9,6}}
  \proofstepwidestar[]{0\neq1}{\rCI{7,10}}
  \proofstepwidestar[]{1\in\mathbb N\land0\neq1}{\rAI{5,11}}
  \proofstepwidestar[]{0\in\mathbb N\land1\in\mathbb N\land0\neq1}{\rAI{1,12}}
\end{tabproofwide}


\par\Needspace{12\baselineskip}
\hypertarget{differences.statement.FDMaxCarrierElements}{}
\FormulaThmDeltaK[Die beiden Elemente des Trägers]{%
  0\in M_0\land1\in M_0
}{FDMaxCarrierElements}{\DeltaRow{Mengen}{M_0\dsep0\dsep1}}
\begin{tabproofwide}
  \proofstepwidestar[]{0\in\{0,1\}}{\FormulaRefAuto{a\in\{a,b\}}[thm]}
  \proofstepwidestar[]{1\in\{0,1\}}{\FormulaRefAuto{b\in\{a,b\}}[thm]}
  \proofstepwidestar[]{0\in M_0}{\FormulaRefAuto{FDMaxCarrier}[def]{1}}
  \proofstepwidestar[]{1\in M_0}{\FormulaRefAuto{FDMaxCarrier}[def]{2}}
  \proofstepwidestar[]{0\in M_0\land1\in M_0}{\rAI{3,4}}
\end{tabproofwide}

\par\Needspace{12\baselineskip}
\hypertarget{differences.statement.FDMaxCarrierSubset}{}
\FormulaThmDeltaK[Der Träger liegt in den natürlichen Zahlen]{%
  M_0\subseteq\mathbb N
}{FDMaxCarrierSubset}{\DeltaRow{Mengen}{M_0}}
\begin{tabproofwide}
  \proofstepwidestar[]{0\in\mathbb N\land1\in\mathbb N\land0\neq1}{\FormulaRefAuto{FDMaxNaturalConstants}[thm]}
  \proofstepwidestar[]{0\in\mathbb N}{\rAEa{1}}
  \proofstepwidestar[]{1\in\mathbb N}{\rAEa{\rAEb{1}}}
  \proofstepwidestar[]{\{0,1\}\subseteq\mathbb N}{\FormulaRefAuto{a\in A,\,b\in A\vdash\{a,b\}\subseteq A}[thm]{2,3}}
  \proofstepwidestar[]{M_0\subseteq\mathbb N}{\FormulaRefAuto{FDMaxCarrier}[def]{4}}
\end{tabproofwide}

\par\Needspace{12\baselineskip}
\hypertarget{differences.statement.FDMaxOrder}{}
\FormulaDefDeltaK[Die natürliche Ordnung auf dem kleinen Träger]{%
  x,y\in M_0\vdash(x\leq_0y\leftrightarrow x\leq_{\mathbb N}y)
}{FDMaxOrder}{%
  \DeltaRow{Mengen}{M_0\dsep x\dsep y}
  \DeltaRow{Relationsoperatoren}{\leq_0\dsep\leq_{\mathbb N}}
  \DeltaRow{Restriktion}{\begin{aligned}[t]
    &\leq_0\text{ ist die Restriktion von }\leq_{\mathbb N}\\[-2pt]
    &\text{auf }M_0
  \end{aligned}}
    [\FormulaRefAuto{OrderRestrictionOperator}[def]]
  \DeltaRow{Teilmenge}{M_0\subseteq\mathbb N}[\FormulaRefAuto{FDMaxCarrierSubset}[thm]]
}

\par\Needspace{12\baselineskip}
\hypertarget{differences.statement.FDMaxTotalOrder}{}
\FormulaThmDeltaK[Die eingeschränkte Ordnung ist total]{%
  \TotOrd{M_0,\leq_0}
}{FDMaxTotalOrder}{%
  \DeltaRow{Mengen}{M_0\dsep x\dsep y}
  \DeltaRow{Relationsoperatoren}{\leq_0\dsep\leq_{\mathbb N}}
}
\begin{tabproofwide}
  \proofstepwidestar[]{\TotOrd{\mathbb N,\leq_{\mathbb N}}}{\FormulaRefAuto{NaturalNumbersTotalOrder}[thm]}
  \proofstepwidestar[]{\PartOrd{\mathbb N,\leq_{\mathbb N}}}{\FormulaRefAuto{Totale Ordnung}[def]{1}}
  \proofstepwidestar[]{M_0\subseteq\mathbb N}{\FormulaRefAuto{FDMaxCarrierSubset}[thm]}
  \proofstepwidestar[]{\PartOrd{M_0,\leq_0}}{\FormulaRefAuto{OrderRestrictionPartialOrder}[thm]{3,2}}
  \proofstepwidestar[5]{x\in M_0}{\rA}
  \proofstepwidestar[6]{y\in M_0}{\rA}
  \proofstepwidestar[5]{x\in\mathbb N}{\FormulaRefAuto{A\subseteq B\dsep x\in A\vdash x\in B}[thm]{3,5}}
  \proofstepwidestar[6]{y\in\mathbb N}{\FormulaRefAuto{A\subseteq B\dsep x\in A\vdash x\in B}[thm]{3,6}}
  \proofstepwidestar[5,6]{x\leq_{\mathbb N}y\lor y\leq_{\mathbb N}x}{\FormulaRefAuto{PeanoOrderComparison}[thm]{7,8}}
  \proofstepwidestar[10]{x\leq_{\mathbb N}y}{\rA}
  \proofstepwidestar[5,6]{x\leq_0y\leftrightarrow x\leq_{\mathbb N}y}{\FormulaRefAuto{FDMaxOrder}[def]{5,6}}
  \proofstepwidestar[5,6,10]{x\leq_0y}{\FormulaRefAuto{P\leftrightarrow Q,Q\vdash P}[thm]{11,10}}
  \proofstepwidestar[5,6,10]{x\leq_0y\lor y\leq_0x}{\rOIa{12}}
  \proofstepwidestar[14]{y\leq_{\mathbb N}x}{\rA}
  \proofstepwidestar[5,6]{y\leq_0x\leftrightarrow y\leq_{\mathbb N}x}{\FormulaRefAuto{FDMaxOrder}[def]{6,5}}
  \proofstepwidestar[5,6,14]{y\leq_0x}{\FormulaRefAuto{P\leftrightarrow Q,Q\vdash P}[thm]{15,14}}
  \proofstepwidestar[5,6,14]{x\leq_0y\lor y\leq_0x}{\rOIb{16}}
  \proofstepwidestar[5,6]{x\leq_0y\lor y\leq_0x}{\rOE{9,10,13,14,17}}
  \proofstepwidestar[5]{y\in M_0\rightarrow(x\leq_0y\lor y\leq_0x)}{\rRI{6,18}}
  \proofstepwidestar[5]{\forall y\in M_0\,(x\leq_0y\lor y\leq_0x)}{\rUI{19}}
  \proofstepwidestar[]{x\in M_0\rightarrow\forall y\in M_0\,(x\leq_0y\lor y\leq_0x)}{\rRI{5,20}}
  \proofstepwidestar[]{\forall x\in M_0\,\forall y\in M_0\,(x\leq_0y\lor y\leq_0x)}{\rUI{21}}
  \proofstepwidestar[]{\TotOrd{M_0,\leq_0}}{\FormulaRefAuto{Totale Ordnung}[def]{4,22}}
\end{tabproofwide}

\par\Needspace{12\baselineskip}
\hypertarget{differences.statement.FDMaxZeroLeast}{}
\FormulaThmDeltaK[Null liegt unter jedem Element]{%
  x\in M_0\vdash0\leq_0x
}{FDMaxZeroLeast}{\DeltaRow{Mengen}{M_0\dsep x}}
\begin{tabproofwide}
  \proofstepwidestar[1]{x\in M_0}{\rA}
  \proofstepwidestar[]{M_0\subseteq\mathbb N}{\FormulaRefAuto{FDMaxCarrierSubset}[thm]}
  \proofstepwidestar[1]{x\in\mathbb N}{\FormulaRefAuto{A\subseteq B\dsep x\in A\vdash x\in B}[thm]{2,1}}
  \proofstepwidestar[1]{0\leq_{\mathbb N}x}{\FormulaRefAuto{PeanoZeroLeq}[thm]{3}}
  \proofstepwidestar[]{0\in M_0\land1\in M_0}{\FormulaRefAuto{FDMaxCarrierElements}[thm]}
  \proofstepwidestar[1]{0\leq_0x\leftrightarrow0\leq_{\mathbb N}x}{\FormulaRefAuto{FDMaxOrder}[def]{\rAEa{5},1}}
  \proofstepwidestar[1]{0\leq_0x}{\FormulaRefAuto{P\leftrightarrow Q,Q\vdash P}[thm]{6,4}}
\end{tabproofwide}

\par\Needspace{16\baselineskip}
\section{Maximum als wohldefinierte Verknüpfung}

\par\Needspace{12\baselineskip}
\hypertarget{differences.statement.FDMaxOperation}{}
\FormulaDefDeltaK[Die Maximumsverknüpfung]{%
  \star_0\coloneqq\max_{M_0,\leq_0}
}{FDMaxOperation}{%
  \DeltaRow{Mengen}{M_0}
  \DeltaRow{Funktionen}{\star_0\dsep\max_{M_0,\leq_0}}
  \DeltaRow{Auswertung}{x,y\in M_0\vdash x\star_0y=\max_{M_0,\leq_0}(x,y)}
    [\FormulaRefAuto{TotalOrderBinaryMaximum}[def]]
}

Die Operation ist damit genau das gewöhnliche Maximum von Null und Eins:
\(\leq_0\) vergleicht die Elemente nach der natürlichen Ordnung. Der
folgende Funktionstyp wird aus der in Band~15 konstruierten binären
Maximumsfunktion gewonnen.

\par\Needspace{12\baselineskip}
\hypertarget{differences.statement.FDMaxFunctionType}{}
\FormulaThmDeltaK[Der Funktionstyp der Maximumsverknüpfung]{%
  \star_0\colon M_0\times M_0\to M_0
}{FDMaxFunctionType}{\DeltaRow{Mengen}{M_0}\DeltaRow{Funktionen}{\star_0}}
\begin{tabproofwide}
  \proofstepwidestar[]{\TotOrd{M_0,\leq_0}}{\FormulaRefAuto{FDMaxTotalOrder}[thm]}
  \proofstepwidestar[]{\max_{M_0,\leq_0}\colon M_0\times M_0\to M_0}{\FormulaRefAuto{TotalOrderBinaryMaximumFunctionType}[thm]{1}}
  \proofstepwidestar[]{\star_0\colon M_0\times M_0\to M_0}{\FormulaRefAuto{FDMaxOperation}[def]{2}}
\end{tabproofwide}

\par\Needspace{12\baselineskip}
\hypertarget{differences.statement.FDMaxClosed}{}
\FormulaThmDeltaK[Abschluss unter dem Maximum]{%
  x\in M_0\dsep y\in M_0\vdash x\star_0y\in M_0
}{FDMaxClosed}{\DeltaRow{Mengen}{M_0\dsep x\dsep y}}
\begin{tabproofwide}
  \proofstepwidestar[1]{x\in M_0}{\rA}
  \proofstepwidestar[2]{y\in M_0}{\rA}
  \proofstepwidestar[]{\TotOrd{M_0,\leq_0}}{\FormulaRefAuto{FDMaxTotalOrder}[thm]}
  \proofstepwidestar[1,2]{x,y\in M_0}{\rAI{1,2}}
  \proofstepwidestar[1,2]{\max_{M_0,\leq_0}(x,y)\in M_0}{\FormulaRefAuto{TotalOrderBinaryMaximumInCarrier}[thm]{3,4}}
  \proofstepwidestar[1,2]{x\star_0y\in M_0}{\FormulaRefAuto{FDMaxOperation}[def]{5}}
\end{tabproofwide}

\par\Needspace{12\baselineskip}
\hypertarget{differences.statement.FDMaxCommutative}{}
\FormulaThmDeltaK[Kommutativität des Maximums]{%
  x\in M_0\dsep y\in M_0\vdash x\star_0y=y\star_0x
}{FDMaxCommutative}{\DeltaRow{Mengen}{M_0\dsep x\dsep y}}
\begin{tabproofwide}
  \proofstepwidestar[1]{x\in M_0}{\rA}
  \proofstepwidestar[2]{y\in M_0}{\rA}
  \proofstepwidestar[]{\TotOrd{M_0,\leq_0}}{\FormulaRefAuto{FDMaxTotalOrder}[thm]}
  \proofstepwidestar[1,2]{x,y\in M_0}{\rAI{1,2}}
  \proofstepwidestar[1,2]{\max_{M_0,\leq_0}(x,y)=\max_{M_0,\leq_0}(y,x)}{\FormulaRefAuto{TotalOrderBinaryMaximumCommutative}[thm]{3,4}}
  \proofstepwidestar[1,2]{x\star_0y=y\star_0x}{\FormulaRefAuto{FDMaxOperation}[def]{5}}
\end{tabproofwide}

\par\Needspace{12\baselineskip}
\hypertarget{differences.statement.FDMaxAssociative}{}
\FormulaThmDeltaK[Assoziativität des Maximums]{%
  \begin{aligned}[t]
    &x\in M_0\dsep y\in M_0\dsep z\in M_0\vdash{}\\[-2pt]
    &(x\star_0y)\star_0z=x\star_0(y\star_0z)
  \end{aligned}
}{FDMaxAssociative}{\DeltaRow{Mengen}{M_0\dsep x\dsep y\dsep z}}
\begin{tabproofwide}
  \proofstepwidestar[1]{x\in M_0}{\rA}
  \proofstepwidestar[2]{y\in M_0}{\rA}
  \proofstepwidestar[3]{z\in M_0}{\rA}
  \proofstepwidestar[]{\TotOrd{M_0,\leq_0}}{\FormulaRefAuto{FDMaxTotalOrder}[thm]}
  \proofstepwidestar[1,2,3]{x,y,z\in M_0}{\rAI{1,\rAI{2,3}}}
  \proofstepwidestarbreak[1,2,3]{%
    \max_{M_0,\leq_0}\bigl(\max_{M_0,\leq_0}(x,y),z\bigr)
      =\max_{M_0,\leq_0}\bigl(x,\max_{M_0,\leq_0}(y,z)\bigr)}{%
    \FormulaRefAuto{TotalOrderBinaryMaximumAssociative}[thm]{4,5}}
  \proofstepwidestar[1,2,3]{(x\star_0y)\star_0z=x\star_0(y\star_0z)}{\FormulaRefAuto{FDMaxOperation}[def]{6}}
\end{tabproofwide}

\par\Needspace{16\baselineskip}
\section{Null als neutrales Element}

\par\Needspace{12\baselineskip}
\hypertarget{differences.statement.FDMaxLeftIdentity}{}
\FormulaThmDeltaK[Null ist linksneutral]{%
  x\in M_0\vdash0\star_0x=x
}{FDMaxLeftIdentity}{\DeltaRow{Mengen}{M_0\dsep x\dsep s\dsep u}}
\begin{tabproofwide}
  \proofstepwidestar[1]{x\in M_0}{\rA}
  \proofstepwidestar[]{\TotOrd{M_0,\leq_0}}{\FormulaRefAuto{FDMaxTotalOrder}[thm]}
  \proofstepwidestar[]{\PartOrd{M_0,\leq_0}}{\FormulaRefAuto{Totale Ordnung}[def]{2}}
  \proofstepwidestar[]{0\in M_0\land1\in M_0}{\FormulaRefAuto{FDMaxCarrierElements}[thm]}
  \proofstepwidestar[]{0\in M_0}{\rAEa{4}}
  \proofstepwidestar[1]{0,x\in M_0}{\rAI{5,1}}
  \proofstepwidestar[1]{\operatorname{SupEx}_{M_0,\leq_0}(0,x)}{\FormulaRefAuto{TotalOrderPairSupremumExistence}[thm]{2,6}}
  \proofstepwidestarbreak[1]{%
    (0,x\in M_0)\land\exists s\in\UB_{M_0,\leq_0}(\{0,x\})\,
      \forall u\in\UB_{M_0,\leq_0}(\{0,x\})\,(s\leq_0u)}{%
    \FormulaRefAuto{PairSupremumExistence}[def]{3,7}}
  \proofstepwidestarbreak[1]{%
    \exists s\in\UB_{M_0,\leq_0}(\{0,x\})\,
      \forall u\in\UB_{M_0,\leq_0}(\{0,x\})\,(s\leq_0u)}{\rAEb{8}}
  \proofstepwidestar[1]{x\leq_0\sup_{M_0,\leq_0}(\{0,x\})}{\FormulaRefAuto{PairSupremumSecondOperand}[thm]{3,5,1,9}}
  \proofstepwidestar[1]{0\leq_0x}{\FormulaRefAuto{FDMaxZeroLeast}[thm]{1}}
  \proofstepwidestar[1]{x\leq_0x}{\FormulaRefAuto{EqRelReflexivity}[ax]{3,1}}
  \proofstepwidestar[1]{\sup_{M_0,\leq_0}(\{0,x\})\leq_0x}{\FormulaRefAuto{PairSupremumUniversal}[thm]{3,5,1,1,9,11,12}}
  \proofstepwidestar[1]{\sup_{M_0,\leq_0}(\{0,x\})=x}{\FormulaRefAuto{OrderAntisymmetry}[ax]{3,13,10}}
  \proofstepwidestar[1]{\max_{M_0,\leq_0}(0,x)=\sup_{M_0,\leq_0}(\{0,x\})}{\FormulaRefAuto{TotalOrderBinaryMaximumAsSupremum}[thm]{2,6}}
  \proofstepwidestar[1]{\max_{M_0,\leq_0}(0,x)=x}{\FormulaRefAuto{a=b\dsep b=c\vdash a=c}[thm]{15,14}}
  \proofstepwidestar[1]{0\star_0x=x}{\FormulaRefAuto{FDMaxOperation}[def]{16}}
\end{tabproofwide}

\par\Needspace{12\baselineskip}
\hypertarget{differences.statement.FDMaxRightIdentity}{}
\FormulaThmDeltaK[Null ist rechtsneutral]{%
  x\in M_0\vdash x\star_0 0=x
}{FDMaxRightIdentity}{\DeltaRow{Mengen}{M_0\dsep x}}
\begin{tabproofwide}
  \proofstepwidestar[1]{x\in M_0}{\rA}
  \proofstepwidestar[]{0\in M_0\land1\in M_0}{\FormulaRefAuto{FDMaxCarrierElements}[thm]}
  \proofstepwidestar[]{0\in M_0}{\rAEa{2}}
  \proofstepwidestar[1]{x\star_0 0=0\star_0x}{\FormulaRefAuto{FDMaxCommutative}[thm]{1,3}}
  \proofstepwidestar[1]{0\star_0x=x}{\FormulaRefAuto{FDMaxLeftIdentity}[thm]{1}}
  \proofstepwidestar[1]{x\star_0 0=x}{\FormulaRefAuto{a=b\dsep b=c\vdash a=c}[thm]{4,5}}
\end{tabproofwide}

\par\Needspace{16\baselineskip}
\section{Die nachgewiesene Monoidstruktur}

Die folgenden Tabellen fassen die punktweisen Gesetze mit ausdrücklich
entladenen Elementannahmen zusammen. Die Strukturbezeichnungen aus den
Bänden~28 und~38 werden erst danach angewendet.

\par\Needspace{12\baselineskip}
\hypertarget{differences.statement.FDMaxSemigroup}{}
\FormulaThmDeltaK[Das Maximum bildet eine Halbgruppe]{%
  \SemigroupAlg{M_0}{\star_0}
}{FDMaxSemigroup}{\DeltaRow{Mengen}{M_0\dsep x\dsep y\dsep z}}
\begin{tabproofwide}
  \proofstepwidestar[]{\star_0\colon M_0\times M_0\to M_0}{\FormulaRefAuto{FDMaxFunctionType}[thm]}
  \proofstepwidestar[2]{x\in M_0}{\rA}
  \proofstepwidestar[3]{y\in M_0}{\rA}
  \proofstepwidestar[2,3]{x\star_0y\in M_0}{\FormulaRefAuto{FDMaxClosed}[thm]{2,3}}
  \proofstepwidestar[2]{y\in M_0\rightarrow x\star_0y\in M_0}{\rRI{3,4}}
  \proofstepwidestar[2]{\forall y\in M_0\,(x\star_0y\in M_0)}{\rUI{5}}
  \proofstepwidestar[]{x\in M_0\rightarrow\forall y\in M_0\,(x\star_0y\in M_0)}{\rRI{2,6}}
  \proofstepwidestar[]{\forall x\in M_0\,\forall y\in M_0\,(x\star_0y\in M_0)}{\rUI{7}}
  \proofstepwidestar[9]{x\in M_0}{\rA}
  \proofstepwidestar[10]{y\in M_0}{\rA}
  \proofstepwidestar[11]{z\in M_0}{\rA}
  \proofstepwidestar[9,10,11]{(x\star_0y)\star_0z=x\star_0(y\star_0z)}{\FormulaRefAuto{FDMaxAssociative}[thm]{9,10,11}}
  \proofstepwidestar[9,10]{z\in M_0\rightarrow(x\star_0y)\star_0z=x\star_0(y\star_0z)}{\rRI{11,12}}
  \proofstepwidestar[9,10]{\forall z\in M_0\,((x\star_0y)\star_0z=x\star_0(y\star_0z))}{\rUI{13}}
  \proofstepwidestarbreak[9]{y\in M_0\rightarrow\forall z\in M_0\,((x\star_0y)\star_0z=x\star_0(y\star_0z))}{\rRI{10,14}}
  \proofstepwidestarbreak[9]{\forall y\in M_0\,\forall z\in M_0\,((x\star_0y)\star_0z=x\star_0(y\star_0z))}{\rUI{15}}
  \proofstepwidestarbreak[]{x\in M_0\rightarrow\forall y\in M_0\,\forall z\in M_0\,((x\star_0y)\star_0z=x\star_0(y\star_0z))}{\rRI{9,16}}
  \proofstepwidestarbreak[]{\forall x\in M_0\,\forall y\in M_0\,\forall z\in M_0\,((x\star_0y)\star_0z=x\star_0(y\star_0z))}{\rUI{17}}
  \proofstepwidestar[]{\SemigroupAlg{M_0}{\star_0}}{\FormulaRefAuto{SemigroupDef}[def]{1,8,18}}
\end{tabproofwide}

\par\Needspace{12\baselineskip}
\hypertarget{differences.statement.FDMaxMonoid}{}
\FormulaThmDeltaK[Das Maximum mit Null bildet ein Monoid]{%
  \MonoidAlg{M_0}{\star_0}{0}
}{FDMaxMonoid}{\DeltaRow{Mengen}{M_0\dsep x}}
\begin{tabproofwide}
  \proofstepwidestar[]{\SemigroupAlg{M_0}{\star_0}}{\FormulaRefAuto{FDMaxSemigroup}[thm]}
  \proofstepwidestar[]{0\in M_0\land1\in M_0}{\FormulaRefAuto{FDMaxCarrierElements}[thm]}
  \proofstepwidestar[]{0\in M_0}{\rAEa{2}}
  \proofstepwidestar[4]{x\in M_0}{\rA}
  \proofstepwidestar[4]{0\star_0x=x}{\FormulaRefAuto{FDMaxLeftIdentity}[thm]{4}}
  \proofstepwidestar[]{x\in M_0\rightarrow0\star_0x=x}{\rRI{4,5}}
  \proofstepwidestar[]{\forall x\in M_0\,(0\star_0x=x)}{\rUI{6}}
  \proofstepwidestar[8]{x\in M_0}{\rA}
  \proofstepwidestar[8]{x\star_0 0=x}{\FormulaRefAuto{FDMaxRightIdentity}[thm]{8}}
  \proofstepwidestar[]{x\in M_0\rightarrow x\star_0 0=x}{\rRI{8,9}}
  \proofstepwidestar[]{\forall x\in M_0\,(x\star_0 0=x)}{\rUI{10}}
  \proofstepwidestar[]{\MonoidAlg{M_0}{\star_0}{0}}{\FormulaRefAuto{MonoidDef}[def]{1,3,7,11}}
\end{tabproofwide}

\par\Needspace{12\baselineskip}
\hypertarget{differences.statement.FDMaxCommutativeMonoid}{}
\FormulaThmDeltaK[Das Maximumsmonoid ist kommutativ]{%
  \CommutativeMonoidAlg{M_0}{\star_0}{0}
}{FDMaxCommutativeMonoid}{\DeltaRow{Mengen}{M_0\dsep x\dsep y}}
\begin{tabproofwide}
  \proofstepwidestar[]{\MonoidAlg{M_0}{\star_0}{0}}{\FormulaRefAuto{FDMaxMonoid}[thm]}
  \proofstepwidestar[2]{x\in M_0}{\rA}
  \proofstepwidestar[3]{y\in M_0}{\rA}
  \proofstepwidestar[2,3]{x\star_0y=y\star_0x}{\FormulaRefAuto{FDMaxCommutative}[thm]{2,3}}
  \proofstepwidestar[2]{y\in M_0\rightarrow x\star_0y=y\star_0x}{\rRI{3,4}}
  \proofstepwidestar[2]{\forall y\in M_0\,(x\star_0y=y\star_0x)}{\rUI{5}}
  \proofstepwidestar[]{x\in M_0\rightarrow\forall y\in M_0\,(x\star_0y=y\star_0x)}{\rRI{2,6}}
  \proofstepwidestar[]{\forall x\in M_0\,\forall y\in M_0\,(x\star_0y=y\star_0x)}{\rUI{7}}
  \proofstepwidestar[]{\CommutativeMonoidAlg{M_0}{\star_0}{0}}{\FormulaRefAuto{CommutativeMonoidDef}[def]{1,8}}
\end{tabproofwide}

\par\Needspace{16\baselineskip}
\section{Die vier Werte und die fehlende Kürzbarkeit}

\par\Needspace{12\baselineskip}
\hypertarget{differences.statement.FDMaxZeroZero}{}
\FormulaThmDeltaK[Maximum von Null und Null]{0\star_0 0=0}{FDMaxZeroZero}{\DeltaRow{Mengen}{M_0}}
\begin{tabproofwide}
  \proofstepwidestar[]{0\in M_0\land1\in M_0}{\FormulaRefAuto{FDMaxCarrierElements}[thm]}
  \proofstepwidestar[]{0\in M_0}{\rAEa{1}}
  \proofstepwidestar[]{0\star_0 0=0}{\FormulaRefAuto{FDMaxLeftIdentity}[thm]{2}}
\end{tabproofwide}

\par\Needspace{12\baselineskip}
\hypertarget{differences.statement.FDMaxZeroOne}{}
\FormulaThmDeltaK[Maximum von Null und Eins]{0\star_0 1=1}{FDMaxZeroOne}{\DeltaRow{Mengen}{M_0}}
\begin{tabproofwide}
  \proofstepwidestar[]{0\in M_0\land1\in M_0}{\FormulaRefAuto{FDMaxCarrierElements}[thm]}
  \proofstepwidestar[]{1\in M_0}{\rAEb{1}}
  \proofstepwidestar[]{0\star_0 1=1}{\FormulaRefAuto{FDMaxLeftIdentity}[thm]{2}}
\end{tabproofwide}

\par\Needspace{12\baselineskip}
\hypertarget{differences.statement.FDMaxOneZero}{}
\FormulaThmDeltaK[Maximum von Eins und Null]{1\star_0 0=1}{FDMaxOneZero}{\DeltaRow{Mengen}{M_0}}
\begin{tabproofwide}
  \proofstepwidestar[]{0\in M_0\land1\in M_0}{\FormulaRefAuto{FDMaxCarrierElements}[thm]}
  \proofstepwidestar[]{1\in M_0}{\rAEb{1}}
  \proofstepwidestar[]{1\star_0 0=1}{\FormulaRefAuto{FDMaxRightIdentity}[thm]{2}}
\end{tabproofwide}

\par\Needspace{12\baselineskip}
\hypertarget{differences.statement.FDMaxOneOne}{}
\FormulaThmDeltaK[Maximum von Eins und Eins]{1\star_0 1=1}{FDMaxOneOne}{\DeltaRow{Mengen}{M_0}}
\begin{tabproofwide}
  \proofstepwidestar[]{0\in M_0\land1\in M_0}{\FormulaRefAuto{FDMaxCarrierElements}[thm]}
  \proofstepwidestar[]{1\in M_0}{\rAEb{1}}
  \proofstepwidestar[]{\TotOrd{M_0,\leq_0}}{\FormulaRefAuto{FDMaxTotalOrder}[thm]}
  \proofstepwidestar[]{\max_{M_0,\leq_0}(1,1)=1}{\FormulaRefAuto{TotalOrderBinaryMaximumIdempotent}[thm]{3,2}}
  \proofstepwidestar[]{1\star_0 1=1}{\FormulaRefAuto{FDMaxOperation}[def]{4}}
\end{tabproofwide}

\par\Needspace{12\baselineskip}
\hypertarget{differences.statement.FDMaxCancellationEquality}{}
\FormulaThmDeltaK[Gleiche Produkte mit rechtem Faktor Eins]{%
  0\star_0 1=1\star_0 1
}{FDMaxCancellationEquality}{\DeltaRow{Mengen}{M_0}}
\begin{tabproofwide}
  \proofstepwidestar[]{0\star_0 1=1}{\FormulaRefAuto{FDMaxZeroOne}[thm]}
  \proofstepwidestar[]{1\star_0 1=1}{\FormulaRefAuto{FDMaxOneOne}[thm]}
  \proofstepwidestar[]{1=1\star_0 1}{\FormulaRefAuto{a=b\vdash b=a}[thm]{2}}
  \proofstepwidestar[]{0\star_0 1=1\star_0 1}{\FormulaRefAuto{a=b\dsep b=c\vdash a=c}[thm]{1,3}}
\end{tabproofwide}

\par\Needspace{12\baselineskip}
\hypertarget{differences.statement.FDMaxNotRightCancellative}{}
\FormulaThmDeltaK[Das Maximum ist nicht rechtskürzbar]{%
  \neg\mathsf{RK\ddot{u}rz}(M_0,\star_0)
}{FDMaxNotRightCancellative}{\DeltaRow{Mengen}{M_0}}
\begin{tabproofwide}
  \proofstepwidestar[1]{\mathsf{RK\ddot{u}rz}(M_0,\star_0)}{\rA}
  \proofstepwidestar[]{0\in M_0\land1\in M_0}{\FormulaRefAuto{FDMaxCarrierElements}[thm]}
  \proofstepwidestar[]{0\in M_0}{\rAEa{2}}
  \proofstepwidestar[]{1\in M_0}{\rAEb{2}}
  \proofstepwidestar[]{0\star_0 1=1\star_0 1}{\FormulaRefAuto{FDMaxCancellationEquality}[thm]}
  \proofstepwidestar[1]{0=1}{\FormulaRefAuto{RightCancellationAxiom}[ax]{1,4,3,4,5}}
  \proofstepwidestar[]{0\in\mathbb N\land1\in\mathbb N\land0\neq1}{\FormulaRefAuto{FDMaxNaturalConstants}[thm]}
  \proofstepwidestar[]{0\neq1}{\rAEb{\rAEb{7}}}
  \proofstepwidestar[1]{\bot}{\rBI{6,8}}
  \proofstepwidestar[]{\neg\mathsf{RK\ddot{u}rz}(M_0,\star_0)}{\rCI{1,9}}
\end{tabproofwide}

\par\Needspace{12\baselineskip}
\hypertarget{differences.statement.FDMaxNotCancellative}{}
\FormulaThmDeltaK[Das Maximum ist nicht kürzbar]{%
  \neg\mathsf{K\ddot{u}rz}(M_0,\star_0)
}{FDMaxNotCancellative}{\DeltaRow{Mengen}{M_0}}
\begin{tabproofwide}
  \proofstepwidestar[1]{\mathsf{K\ddot{u}rz}(M_0,\star_0)}{\rA}
  \proofstepwidestar[1]{\mathsf{RK\ddot{u}rz}(M_0,\star_0)}{\FormulaRefAuto{CancellativeRightAxiom}[ax]{1}}
  \proofstepwidestar[]{\neg\mathsf{RK\ddot{u}rz}(M_0,\star_0)}{\FormulaRefAuto{FDMaxNotRightCancellative}[thm]}
  \proofstepwidestar[1]{\bot}{\rBI{2,3}}
  \proofstepwidestar[]{\neg\mathsf{K\ddot{u}rz}(M_0,\star_0)}{\rCI{1,4}}
\end{tabproofwide}

\par\Needspace{16\baselineskip}
\section{Drei konkrete Kreuzsummenvergleiche}

\par\Needspace{12\baselineskip}
\hypertarget{differences.statement.FDMaxPairCarrier}{}
\FormulaDefDeltaK[Der Träger der Paare im Gegenbeispiel]{%
  D_0\coloneqq M_0\times M_0
}{FDMaxPairCarrier}{\DeltaRow{Mengen}{D_0\dsep M_0}}

\par\Needspace{12\baselineskip}
\hypertarget{differences.statement.FDMaxRelation}{}
\FormulaDefDeltaK[Abkürzung für die Kreuzsummenrelation]{%
  \begin{aligned}[t]
    &p,q\in D_0\vdash{}\\[-2pt]
    &p\rho_0q\coloneqq
      p\FormalDifferenceRel{M_0,\star_0,0}q
  \end{aligned}
}{FDMaxRelation}{%
  \DeltaRow{Mengen}{D_0\dsep M_0\dsep p\dsep q}
  \DeltaRow{Relationsoperatoren}{\rho_0\dsep\FormalDifferenceRel{M_0,\star_0,0}}
  \DeltaRow{Ausgangsstruktur}{\CommutativeMonoidAlg{M_0}{\star_0}{0}}
    [\FormulaRefAuto{FDMaxCommutativeMonoid}[thm]]
  \DeltaRow{Paarträger}{\FormalDifferenceCarrier{M_0,\star_0,0}=M_0\times M_0}
    [\FormulaRefAuto{GroupCompletionPairCarrier}[def]]
}

\par\Needspace{12\baselineskip}
\hypertarget{differences.statement.FDMaxPairTyped}{}
\FormulaThmDeltaK[Paare von Trägerelementen gehören zum Paarträger]{%
  a\in M_0\dsep b\in M_0\vdash(a,b)\in D_0
}{FDMaxPairTyped}{\DeltaRow{Mengen}{M_0\dsep D_0\dsep a\dsep b}}
\begin{tabproofwide}
  \proofstepwidestar[1]{a\in M_0}{\rA}
  \proofstepwidestar[2]{b\in M_0}{\rA}
  \proofstepwidestar[1,2]{a\in M_0\land b\in M_0}{\rAI{1,2}}
  \proofstepwidestar[1,2]{(a,b)\in M_0\times M_0}{\FormulaRefAuto{(a,b)\in A\times B\eqvdash a\in A\land b\in B}[thm]{3}}
  \proofstepwidestar[1,2]{(a,b)\in D_0}{\FormulaRefAuto{FDMaxPairCarrier}[def]{4}}
\end{tabproofwide}

\par\Needspace{12\baselineskip}
\hypertarget{differences.statement.FDMaxRelationChar}{}
\FormulaThmDeltaK[Der konkrete Kreuzsummenvergleich]{%
  \begin{aligned}[t]
    &a\in M_0\dsep b\in M_0\dsep c\in M_0\dsep d\in M_0\vdash{}\\[-2pt]
    &(a,b)\rho_0(c,d)\leftrightarrow a\star_0d=b\star_0c
  \end{aligned}
}{FDMaxRelationChar}{\DeltaRow{Mengen}{M_0\dsep D_0\dsep a\dsep b\dsep c\dsep d}}
\begin{tabproofwide}
  \proofstepwidestar[1]{a\in M_0}{\rA}
  \proofstepwidestar[2]{b\in M_0}{\rA}
  \proofstepwidestar[3]{c\in M_0}{\rA}
  \proofstepwidestar[4]{d\in M_0}{\rA}
  \proofstepwidestar[]{\CommutativeMonoidAlg{M_0}{\star_0}{0}}{\FormulaRefAuto{FDMaxCommutativeMonoid}[thm]}
  \proofstepwidestar[1,2,3,4]{a,b,c,d\in M_0}{\rAI{1,\rAI{2,\rAI{3,4}}}}
  \proofstepwidestarbreak[1,2,3,4]{%
    (a,b)\FormalDifferenceRel{M_0,\star_0,0}(c,d)
      \leftrightarrow a\star_0d=b\star_0c}{%
    \FormulaRefAuto{GroupCompletionPairRelationChar}[thm]{5,6}}
  \proofstepwidestar[1,2]{(a,b)\in D_0}{\FormulaRefAuto{FDMaxPairTyped}[thm]{1,2}}
  \proofstepwidestar[3,4]{(c,d)\in D_0}{\FormulaRefAuto{FDMaxPairTyped}[thm]{3,4}}
  \proofstepwidestar[1,2,3,4]{(a,b)\rho_0(c,d)\leftrightarrow a\star_0d=b\star_0c}{\FormulaRefAuto{FDMaxRelation}[def]{8,9,7}}
\end{tabproofwide}

\par\Needspace{12\baselineskip}
\hypertarget{differences.statement.FDMaxFirstRelated}{}
\FormulaThmDeltaK[Die Paare Null--Null und Eins--Eins sind verwandt]{%
  (0,0)\rho_0(1,1)
}{FDMaxFirstRelated}{\DeltaRow{Mengen}{M_0\dsep D_0}}
\begin{tabproofwide}
  \proofstepwidestar[]{0\in M_0\land1\in M_0}{\FormulaRefAuto{FDMaxCarrierElements}[thm]}
  \proofstepwidestar[]{0\in M_0}{\rAEa{1}}
  \proofstepwidestar[]{1\in M_0}{\rAEb{1}}
  \proofstepwidestar[]{(0,0)\rho_0(1,1)\leftrightarrow0\star_0 1=0\star_0 1}{\FormulaRefAuto{FDMaxRelationChar}[thm]{2,2,3,3}}
  \proofstepwidestar[]{0\star_0 1=0\star_0 1}{\rII}
  \proofstepwidestar[]{(0,0)\rho_0(1,1)}{\FormulaRefAuto{P\leftrightarrow Q,Q\vdash P}[thm]{4,5}}
\end{tabproofwide}

\par\Needspace{12\baselineskip}
\hypertarget{differences.statement.FDMaxSecondRelated}{}
\FormulaThmDeltaK[Die Paare Eins--Eins und Null--Eins sind verwandt]{%
  (1,1)\rho_0(0,1)
}{FDMaxSecondRelated}{\DeltaRow{Mengen}{M_0\dsep D_0}}
\begin{tabproofwide}
  \proofstepwidestar[]{0\in M_0\land1\in M_0}{\FormulaRefAuto{FDMaxCarrierElements}[thm]}
  \proofstepwidestar[]{0\in M_0}{\rAEa{1}}
  \proofstepwidestar[]{1\in M_0}{\rAEb{1}}
  \proofstepwidestar[]{(1,1)\rho_0(0,1)\leftrightarrow1\star_0 1=1\star_0 0}{\FormulaRefAuto{FDMaxRelationChar}[thm]{3,3,2,3}}
  \proofstepwidestar[]{1\star_0 1=1}{\FormulaRefAuto{FDMaxOneOne}[thm]}
  \proofstepwidestar[]{1\star_0 0=1}{\FormulaRefAuto{FDMaxOneZero}[thm]}
  \proofstepwidestar[]{1=1\star_0 0}{\FormulaRefAuto{a=b\vdash b=a}[thm]{6}}
  \proofstepwidestar[]{1\star_0 1=1\star_0 0}{\FormulaRefAuto{a=b\dsep b=c\vdash a=c}[thm]{5,7}}
  \proofstepwidestar[]{(1,1)\rho_0(0,1)}{\FormulaRefAuto{P\leftrightarrow Q,Q\vdash P}[thm]{4,8}}
\end{tabproofwide}

\par\Needspace{12\baselineskip}
\hypertarget{differences.statement.FDMaxEndpointsNotRelated}{}
\FormulaThmDeltaK[Die Paare Null--Null und Null--Eins sind nicht verwandt]{%
  \neg\bigl((0,0)\rho_0(0,1)\bigr)
}{FDMaxEndpointsNotRelated}{\DeltaRow{Mengen}{M_0\dsep D_0}}
\begin{tabproofwide}
  \proofstepwidestar[1]{(0,0)\rho_0(0,1)}{\rA}
  \proofstepwidestar[]{0\in M_0\land1\in M_0}{\FormulaRefAuto{FDMaxCarrierElements}[thm]}
  \proofstepwidestar[]{0\in M_0}{\rAEa{2}}
  \proofstepwidestar[]{1\in M_0}{\rAEb{2}}
  \proofstepwidestar[]{(0,0)\rho_0(0,1)\leftrightarrow0\star_0 1=0\star_0 0}{\FormulaRefAuto{FDMaxRelationChar}[thm]{3,3,3,4}}
  \proofstepwidestar[1]{0\star_0 1=0\star_0 0}{\FormulaRefAuto{P\leftrightarrow Q,P\vdash Q}[thm]{5,1}}
  \proofstepwidestar[]{0\star_0 1=1}{\FormulaRefAuto{FDMaxZeroOne}[thm]}
  \proofstepwidestar[]{0\star_0 0=0}{\FormulaRefAuto{FDMaxZeroZero}[thm]}
  \proofstepwidestar[1]{1=0\star_0 0}{\rIE{7,6}}
  \proofstepwidestar[1]{1=0}{\rIE{8,9}}
  \proofstepwidestar[1]{0=1}{\FormulaRefAuto{a=b\vdash b=a}[thm]{10}}
  \proofstepwidestar[]{0\in\mathbb N\land1\in\mathbb N\land0\neq1}{\FormulaRefAuto{FDMaxNaturalConstants}[thm]}
  \proofstepwidestar[]{0\neq1}{\rAEb{\rAEb{12}}}
  \proofstepwidestar[1]{\bot}{\rBI{11,13}}
  \proofstepwidestar[]{\neg\bigl((0,0)\rho_0(0,1)\bigr)}{\rCI{1,14}}
\end{tabproofwide}

\par\Needspace{16\baselineskip}
\section{Der Transitivitätswiderspruch}

\par\Needspace{12\baselineskip}
\hypertarget{differences.statement.FDMaxNotTransitive}{}
\FormulaThmDeltaK[Die Kreuzsummenrelation ist nicht transitiv]{%
  \begin{aligned}[t]
    &\neg\forall p\in D_0\,\forall q\in D_0\,\forall r\in D_0\,\bigl(\\[-2pt]
    &\qquad(p\rho_0q\land q\rho_0r)\rightarrow p\rho_0r\bigr)
  \end{aligned}
}{FDMaxNotTransitive}{\DeltaRow{Mengen}{D_0\dsep M_0\dsep p\dsep q\dsep r}}
\begin{tabproofwide}
  \proofstepwidestarbreak[1]{%
    \forall p\in D_0\,\forall q\in D_0\,\forall r\in D_0\,
      ((p\rho_0q\land q\rho_0r)\rightarrow p\rho_0r)}{\rA}
  \proofstepwidestar[]{0\in M_0\land1\in M_0}{\FormulaRefAuto{FDMaxCarrierElements}[thm]}
  \proofstepwidestar[]{0\in M_0}{\rAEa{2}}
  \proofstepwidestar[]{1\in M_0}{\rAEb{2}}
  \proofstepwidestar[]{(0,0)\in D_0}{\FormulaRefAuto{FDMaxPairTyped}[thm]{3,3}}
  \proofstepwidestar[]{(1,1)\in D_0}{\FormulaRefAuto{FDMaxPairTyped}[thm]{4,4}}
  \proofstepwidestar[]{(0,1)\in D_0}{\FormulaRefAuto{FDMaxPairTyped}[thm]{3,4}}
  \proofstepwidestarbreak[1]{%
    (0,0)\in D_0\rightarrow\forall q\in D_0\,\forall r\in D_0\,
      (((0,0)\rho_0q\land q\rho_0r)\rightarrow(0,0)\rho_0r)}{\rUE{1}}
  \proofstepwidestarbreak[1]{%
    \forall q\in D_0\,\forall r\in D_0\,
      (((0,0)\rho_0q\land q\rho_0r)\rightarrow(0,0)\rho_0r)}{\rRE{8,5}}
  \proofstepwidestarbreak[1]{%
    (1,1)\in D_0\rightarrow\forall r\in D_0\,
      (((0,0)\rho_0(1,1)\land(1,1)\rho_0r)\rightarrow(0,0)\rho_0r)}{\rUE{9}}
  \proofstepwidestarbreak[1]{%
    \forall r\in D_0\,
      (((0,0)\rho_0(1,1)\land(1,1)\rho_0r)\rightarrow(0,0)\rho_0r)}{\rRE{10,6}}
  \proofstepwidestarbreak[1]{%
    (0,1)\in D_0\rightarrow
      (((0,0)\rho_0(1,1)\land(1,1)\rho_0(0,1))\rightarrow(0,0)\rho_0(0,1))}{\rUE{11}}
  \proofstepwidestarbreak[1]{%
    ((0,0)\rho_0(1,1)\land(1,1)\rho_0(0,1))\rightarrow(0,0)\rho_0(0,1)}{\rRE{12,7}}
  \proofstepwidestar[]{(0,0)\rho_0(1,1)}{\FormulaRefAuto{FDMaxFirstRelated}[thm]}
  \proofstepwidestar[]{(1,1)\rho_0(0,1)}{\FormulaRefAuto{FDMaxSecondRelated}[thm]}
  \proofstepwidestar[]{(0,0)\rho_0(1,1)\land(1,1)\rho_0(0,1)}{\rAI{14,15}}
  \proofstepwidestar[1]{(0,0)\rho_0(0,1)}{\rRE{13,16}}
  \proofstepwidestar[]{\neg\bigl((0,0)\rho_0(0,1)\bigr)}{\FormulaRefAuto{FDMaxEndpointsNotRelated}[thm]}
  \proofstepwidestar[1]{\bot}{\rBI{17,18}}
  \proofstepwidestarbreak[]{%
    \neg\forall p\in D_0\,\forall q\in D_0\,\forall r\in D_0\,
      ((p\rho_0q\land q\rho_0r)\rightarrow p\rho_0r)}{\rCI{1,19}}
\end{tabproofwide}

Damit sind sowohl die vollständige Ausgangsstruktur des Beispiels als auch
die drei in der Lesefassung benutzten Vergleiche durch Tabellen belegt.
Bereits auf diesem Träger scheitert also die Transitivität der unveränderten
Kreuzsummenrelation, sobald die Kürzbarkeit fehlt.
